\section{Flow Matching: transportar el ruido hasta el video a lo largo de un campo de velocidades}

Con un latent manejable, el siguiente paso es definir qué predice la red. Flow Matching (emparejamiento de flujos) formula la generación como un transporte en tiempo continuo: parte del ruido gaussiano $X_0$ y avanza, a lo largo del campo de velocidades que aprende el modelo, hacia el latent de datos $X_1$. Al igual que diffusion, aprende el proceso generativo inverso, pero su objetivo de entrenamiento y la interpretación del muestreo se acercan más a las ecuaciones diferenciales ordinarias.

\subsection{De la intuición de diffusion al campo de velocidades}

La sección anterior ya comprimió los píxeles en un latent modelable; esta sección responde a la pregunta «cómo obtener ese latent a partir del ruido». La figura general que sigue presenta primero la estructura común: el proceso directo conecta los datos con el ruido y la red neuronal aprende el camino inverso. La diferencia principal no es «si hay o no eliminación de ruido en varios pasos», sino cómo se definen el camino intermedio, el objetivo de predicción y el método de integración numérica.

\lecturefigure{slide-187-flow-matching-overview.jpg}{Diffusion y Flow Matching aprenden el proceso inverso que regresa del ruido a los datos}{00:22:06}

El camino lineal más sencillo se escribe como
\[
X_t=(1-t)X_0+tX_1,
\qquad t\sim\mathcal U(0,1),
\]
donde $X_0\sim\mathcal N(0,I)$ es el ruido, $X_1$ es el latent de un video real y $t$ es el tiempo continuo. La velocidad objetivo de este camino es
\[
V_t=\frac{dX_t}{dt}=X_1-X_0.
\]
El camino lineal simplifica el objetivo de velocidad, pero en el entrenamiento real también pueden modificarse el muestreo del tiempo y la parametrización del camino.

\begin{importantbox}{Qué aprende Flow Matching en el entrenamiento}
El modelo no predice directamente el video final ni clasifica el timestep; aprende hacia qué dirección debe moverse en cualquier estado intermedio $X_t$ y bajo cualquier condición $c$. El campo de velocidades condicional $v_\theta(X_t,t,c)$ vincula el texto, el tiempo y el noisy latent actual.
\end{importantbox}

\subsection{Entrenamiento: supervisión local sobre un camino muestreado al azar}

La sección anterior definió el campo de velocidades continuo; esta explica cómo se convierte en una supervisión ejecutable por minibatch. El entrenamiento no necesita simular la trayectoria completa de 0 a 1: en cada paso se muestrean datos, ruido y un instante de tiempo, se construye $X_t$ y se supervisa la velocidad local. La slide y la fórmula siguientes desarrollan «cómo las regresiones locales se integran en un transporte global».

\lecturefigure{slide-188-flow-matching-training.jpg}{Entrenamiento de Flow Matching: datos, ruido, muestreo del tiempo y regresión de la velocity}{00:23:58}

La pérdida típica es
\[
\mathcal L_{\mathrm{FM}}(\theta)
=\mathbb E_{X_1,X_0,t,c}
\left[\left\|v_\theta(X_t,t,c)-(X_1-X_0)\right\|_2^2\right].
\]
$\theta$ son los parámetros del Transformer, $c$ es la condición de texto y la esperanza abarca los datos de entrenamiento, el ruido y el muestreo del tiempo. El modelo solo ve estados locales, pero a partir de la supervisión en todos los instantes de tiempo compone el campo de transporte global.

\begin{lstlisting}[language=Python,caption={Pseudocódigo de un paso de entrenamiento de Flow Matching condicional}]
video, prompt = sample_batch()
x1 = tae.encode(video)
x0 = randn_like(x1)
t = sample_time(batch_size)
xt = (1 - t) * x0 + t * x1
target_velocity = x1 - x0
prediction = transformer(xt, timestep=t, text=encode_text(prompt))
loss = mse(prediction, target_velocity)
loss.backward()
\end{lstlisting}

\subsection{Inferencia: la discretización de la ODE es un presupuesto de latencia}

Lo que se obtiene del entrenamiento es la dirección local en cualquier instante de tiempo; para generar una muestra hay que acumular esas direcciones a lo largo del camino. Esta sección lleva el objetivo continuo a un sistema discreto: qué timesteps elegir, qué solver usar y cuántas veces ejecutar el backbone son decisiones que modifican a la vez el error de la trayectoria, la latencia en tiempo de reloj y el uso de memoria de la GPU.

Una vez entrenado el campo de velocidades, el muestreo parte de $\hat X_0\sim\mathcal N(0,I)$ y resuelve
\[
\frac{d\hat X_t}{dt}=v_\theta(\hat X_t,t,c).
\]
La actualización de Euler de primer orden más sencilla es
\[
\hat X_{t_{k+1}}
=\hat X_{t_k}+\Delta t_k\,v_\theta(\hat X_{t_k},t_k,c),
\qquad \Delta t_k=t_{k+1}-t_k.
\]
$k$ denota el paso discreto de resolución; en general, más pasos dan una mejor aproximación, pero cada paso requiere ejecutar una vez un Transformer grande.

\lecturefigure{slide-189-flow-matching-inference.jpg}{Inferencia de Flow Matching: iterar con un ODE solver desde el ruido hasta el latent de datos}{00:24:41}

\begin{warningbox}{Entrenar en tiempo continuo no equivale a calcular en tiempo continuo}
Durante el entrenamiento el modelo puede muestrear $t$ continuo, pero en la inferencia real todavía hay que elegir un número finito de instantes de tiempo y un solver. Muestrear con pocos pasos reduce la latencia pero aumenta el error de discretización; muestrear con muchos pasos eleva el costo. Al reportar la velocidad de «generar un video» deben indicarse también la resolución, el número de cuadros, el solver y el número de pasos.
\end{warningbox}

\teachervoice{00:32:28--00:33:54, pregunta de seguimiento en la sesión de preguntas sobre los denoising steps. El docente explica que la variable de tiempo es continua durante el entrenamiento y que en la inferencia solo se muestrea un conjunto de puntos discretos; más pasos se acercan más a la trayectoria objetivo, pero aumentan directamente la latency.}

\subsection{Resumen de la sección}

Flow Matching construye una supervisión local de velocidad a partir del camino entre el ruido aleatorio y el latent de datos, y en la inferencia obtiene las muestras integrando con un ODE solver. Que la fórmula sea sencilla no significa que la inferencia sea barata: en cada paso de tiempo hay que ejecutar un backbone del orden de 30B, por lo que el número de pasos discretos es el principal parámetro de ajuste entre calidad y latencia.
